\documentclass[11pt,a4paper]{article}
\usepackage[margin=25mm,headheight=16pt]{geometry}
\usepackage{fontspec,kotex,xcolor,hyperref,fancyhdr}
\setmainfont{DejaVu Serif}
\setmainhangulfont{Noto Sans CJK KR}
\definecolor{navy}{HTML}{173E59}
\pagestyle{fancy}\fancyhf{}
\fancyhead[L]{\small GuardSynth CoC}
\fancyhead[R]{\small 제목·초록 검토안 · 한글}
\fancyfoot[L]{\small v0.1 · 2026-09-30}
\fancyfoot[R]{\small\thepage}
\renewcommand{\headrulewidth}{0.3pt}
\setlength{\parindent}{0pt}
\setlength{\emergencystretch}{3em}
\linespread{1.15}
\hypersetup{pdftitle={GuardSynth: CoC 기반 VLM 학습 보강을 위한 행동 제약의 정형 명세와 검증},pdfsubject={Title and abstract for author review},pdfauthor={GuardSynth project}}
\begin{document}
\vspace*{8mm}
{\Large\bfseries\color{navy}\raggedright GuardSynth: CoC 기반 VLM 학습 보강을 위한 행동 제약의 정형 명세와 검증\par}
\vspace{8mm}
{\large\bfseries 초록\par}
\vspace{3mm}
Alpamayo와 같은 추론 기반 자율주행 모델은 Chain of Causation(CoC)을 통해 장면에 대한 인과적 추론과 행동 예측을 연결한다. 그러나 상황과 행동의 이유를 설명하는 CoC 기반 학습만으로는 올바른 행동 선택과 제약 준수를 충분히 학습하는 데 한계가 있다. 본 논문은 이러한 한계를 보완하기 위해 행동 제약을 체계적으로 명세·검증하고 CoC 학습 데이터에 결합하는 GuardSynth를 제안한다. 행동 제약의 범위·특성·구성 요소를 정의하고, 적용 조건, 금지 행동, 유지·해제 조건, 시간 및 근거를 Evidence-Bound Lifecycle Contracts(EBLC)로 명세한다. 이어 제한된 논리 모델에서 일관성과 지정 속성을 검증하고, 명세에 대응하는 통제 자연어(Controlled Natural Language, CNL)를 생성하여 CoC를 보강한다. 통제된 합성 시각·시간 과제에서 시각언어모델(VLM)인 Qwen3-VL-2B-Instruct를 동일한 LoRA 학습 예산과 세 개의 시드로 비교하였으며, 제약 추가 조건에는 학습과 평가 모두에서 제약을 제공하였다. CoC 단독, 기존 자연어 제약 추가, EBLC 기반 CNL 추가 조건의 평균 행동 선택 정확도는 각각 49.86\%, 96.81\%, 98.06\%였으며, 제약 위반율은 19.17\%, 1.39\%, 0.83\%였다. 결과는 행동 제약의 정형 명세·검증을 CoC 학습에 연결하는 체계와 명시적 제약 제공에 따른 행동 선택 개선을 보여준다.
\end{document}
